\chapter{Background: Standards, Practice, and Perception}
\label{ch:background}

This chapter gathers the material the framework depends on. It is not a survey of film sound studies. It collects, from several fields that rarely cite each other, the pieces needed to state the problem of occupation precisely: how loudness is measured and regulated, how cinema reproduction is calibrated and why home reproduction differs, how film sound is organised in practice, and what psychoacoustics and noise research say about intelligibility, annoyance, and recovery.

\section{Measuring loudness}

The reference algorithm for programme loudness is ITU-R BS.1770 \citep{itu1770}. It applies a frequency weighting (the K-weighting, a high-shelf and high-pass pair approximating the head's acoustic effect and reduced low-frequency sensitivity), computes mean-square power per channel, sums channels with weights that favour the surrounds slightly, and gates out quiet blocks so that silence does not dilute the result. Its output is expressed in LUFS (or LKFS, the same unit under a different name). European broadcast practice under EBU R\,128 normalises programmes to $-23$~LUFS integrated loudness \citep{ebu128}; North American practice under ATSC A/85 uses $-24$~LKFS \citep{atsc85}. In the United States the Commercial Advertisement Loudness Mitigation Act made A/85 compliance for advertisements a legal requirement \citep{calm2010}.

Integrated loudness is a single number for a whole programme. EBU Tech~3342 supplements it with Loudness Range (LRA), a statistic of the distribution of short-term loudness values after gating, computed as the spread between the 10th and 95th percentiles \citep{ebu3342}. LRA answers how widely the loudness varies. It does not answer how that variation is arranged in time. A programme alternating every ten seconds between quiet and loud passages and a programme with one quiet hour followed by one loud hour can have the same integrated loudness and the same LRA. For the purposes of this study that is the decisive limitation of the standard measures: they are distributional, and occupation is a temporal property.

\section{Dialogue as reference}

Broadcast and home formats already acknowledge that dialogue is the natural anchor for listening level. The AC-3 bitstream carries a \emph{dialnorm} parameter describing the average level of dialogue in the programme, which decoders use to bring dialogue to a common reproduction level; the same format carries dynamic range control profiles that compress loud and quiet passages toward that anchor when the listener selects a reduced-range mode \citep{atsc52}. Consumer ``night mode'' settings are descendants of this mechanism. Their existence is itself evidence, of an institutional kind, that the gap between dialogue level and effects level is a recognised home-listening problem rather than an idiosyncratic one.

Speech intelligibility has its own measurement tradition. The Speech Intelligibility Index \citep{ansi1997} weights frequency bands by their importance to understanding; the most heavily weighted bands lie in the region from roughly one to four kilohertz, where consonant cues are carried. A reproduction chain weak in that region reduces intelligibility at a given level, which in turn raises the level a listener must choose to follow speech. This point returns in Chapter~\ref{ch:chain}.

\section{Cinema calibration and the home}

Theatrical reproduction is calibrated. Under SMPTE practice each screen channel is aligned so that a pink-noise signal at a specified digital reference level produces a specified sound pressure level at the reference listening position, conventionally 85~dB(C) for $-20$~dBFS \citep{smpte200}. Mixing stages are aligned the same way. A film mix is therefore authored for a known relationship between digital level and acoustic level. Its dynamic range, including the distance between dialogue and the loudest effects, is chosen with that relationship in mind.

Home reproduction has no such alignment. The listener chooses the level, typically lower than theatrical reference, and chooses it by ear, usually for dialogue. Equal-loudness relationships are level dependent \citep{fastl2007}, so a mix balanced at reference level is heard with a different spectral balance at lower levels. More importantly for this study, the listener sets the anchor, and the anchor is speech. Everything the mix places far above speech arrives far above the listener's chosen level. A dynamic range that is intended in a calibrated theatre becomes, in an uncalibrated room, a sequence of exceedances that the listener manages by hand.\footnote{This is an established fact about calibration practice, not a claim that any particular film was mixed badly. A mix can be entirely correct for its intended room and still be difficult in an uncalibrated one. The monograph distinguishes these throughout.}

\section{Film sound practice and theory}

The film sound literature of the digital era documents the expansion of dynamic range and low-frequency capacity that multichannel formats made available \citep{sergi2004,kerins2011}. The low-frequency effects channel and the headroom of digital formats allow effects to be much louder, relative to dialogue, than optical soundtracks permitted.

Two ideas from practitioners and theorists are central here. The first is Chion's observation that cinema is \emph{vococentric}: its sound is organised around the voice, and the ordinary work of mixing is to keep speech intelligible against everything else \citep{chion1994}. The complaint at the root of this study is, in Chion's terms, a failure of vococentrism at the point of reproduction. The second is Murch's account of density \citep{murch2005}. Murch describes a practical limit on how many simultaneous layers of sound of the same kind an audience can follow as distinct, and a strategy of distributing layers across contrasting kinds (dialogue, music, effects of different spectral and rhythmic character) so that density does not collapse into undifferentiated mass. Murch's concern is clarity within a moment. The present study's concern is continuity across moments. The two are complementary: a mix can be clear in every individual moment and still occupy the listener for too long.

\section{Psychoacoustics}

Classical psychoacoustics offers measures closer to perception than digital level. Fastl and Zwicker's models of loudness (in sones), sharpness (weighting of high-frequency energy), roughness (fast amplitude modulation), and fluctuation strength (slow modulation), and the psychoacoustic annoyance index built from them, are the standard descriptive vocabulary \citep{fastl2007}. The analysis pipeline of Chapter~\ref{ch:pipeline} uses cruder spectral descriptors (centroid, flatness, zero-crossing rate), and Chapter~\ref{ch:failures} shows one consequence of that choice. Replacing them with the psychoacoustic measures is part of the programme proposed in Chapter~\ref{ch:response}.

Two further psychoacoustic facts matter. Forward masking reduces sensitivity to quiet sounds for a short interval following a loud one \citep{moore2012}, so speech arriving immediately after an impact is harder to hear than the same speech after silence. And the auditory system organises sound into streams by onset, spectrum, and continuity \citep{bregman1990}; abrupt onsets are strong cues for the formation of a new stream and draw attention. A soundtrack dense with abrupt onsets is, in stream terms, repeatedly asking for attention.

\section{Noise, annoyance, and control}

Community noise research establishes that annoyance grows with level, with the number of events, and with their duration \citep{kryter1985,miedema2001}. The finding most relevant here is older and experimental. Glass and Singer showed that the aftereffects of noise exposure on subsequent task performance depended strongly on predictability and on perceived control: participants who believed they could stop the noise showed markedly reduced aftereffects even when they did not use the control \citep{glass1972}. A volume control is exactly such a means of control. Its use during a film is therefore not only a response to excessive sound but a measurable act of regaining control over it, which is part of why Chapter~\ref{ch:response} treats volume behaviour as the natural dependent variable.

\section{Sound, affect, and the soundscape}

Two broader literatures frame the problem. Goodman's account of sonic warfare treats low frequencies and sudden sound as instruments of affective control that operate below deliberate attention \citep{goodman2010}. Schafer's distinction between hi-fi soundscapes, in which discrete sounds are heard clearly against a quiet ground, and lo-fi soundscapes, in which signals crowd each other and perspective is lost, gives a vocabulary for what a highly occupied soundtrack does to a listening room \citep{schafer1977}. Finally, the music-industry ``loudness war'' shows how competitive pressure can push a production norm toward sustained high level at the expense of contrast \citep{vickers2010}. Whether an analogous pressure operates in film mixing is outside this study's evidence, but the analogy explains why the question is worth asking.

\section{The gap}

None of these literatures measures the quantity this study is concerned with: the temporal structure of a soundtrack's exceedances over its own dialogue level, and the behaviour that structure forces on a listener reproducing it at an uncalibrated level through a particular chain. Loudness standards measure distributions. Film theory describes intentions and effects. Psychoacoustics and noise research describe perception and annoyance under controlled conditions. The framework in the next chapter is an attempt to connect them.
